\documentclass[12pt]{article}
\usepackage[letterpaper,top=0.75in,bottom=0.75in,left=0.78in,right=1.55in,marginparwidth=1.15in,marginparsep=0.18in]{geometry}
\usepackage{fontspec}
\setmainfont{Latin Modern Roman}
\usepackage{microtype}
\usepackage{fancyhdr}
\usepackage{titlesec}
\usepackage{enumitem}
\usepackage{xcolor}
\usepackage[hidelinks]{hyperref}
\setlength{\parindent}{1.05em}
\setlength{\parskip}{0.32em}
\setlength{\headheight}{15pt}
\titleformat{\section}{\large\bfseries\scshape}{}{0pt}{}
\titleformat{\subsection}{\normalsize\bfseries}{}{0pt}{}
\pagestyle{fancy}
\fancyhf{}
\fancyhead[L]{Whately Logic: Julius Caesar Lit-Anchor}
\fancyhead[R]{Lesson 6}
\fancyfoot[C]{\thepage}
\renewcommand{\headrulewidth}{0.4pt}

\newcommand{\focusbox}[1]{\par\vspace{0.35em}\noindent\begingroup\setlength{\fboxsep}{6pt}\colorbox{gray!12}{\begin{minipage}{\dimexpr\linewidth-2\fboxsep\relax}#1\end{minipage}}\endgroup\par\vspace{0.45em}}
\newcommand{\studentline}{\par\vspace{0.68em}\hrule\vspace{0.72em}}
\newcommand{\shortline}{\par\vspace{0.42em}\hrule\vspace{0.55em}}
\newcommand{\gloss}[1]{\marginpar{\scriptsize\raggedright\setlength{\baselineskip}{7.8pt}\color{gray!70!black}#1}}
\newcommand{\speaker}[1]{\par\vspace{0.35em}\noindent\textsc{#1}\par\vspace{-0.15em}}

\begin{document}
\begin{center}
{\LARGE\bfseries Julius Caesar Lit-Anchor}\par\vspace{0.18em}
{\large\scshape Valid Form, Unproved Premise}\par\vspace{0.35em}
{Lesson 6: Valid Does Not Mean True}\par\vspace{0.55em}
\rule{0.82\textwidth}{0.4pt}
\end{center}

\section*{Purpose}
Lesson 6 asks you to separate two questions that are often confused: Does the conclusion follow? Are the premises true? This lit-anchor returns to \textit{Julius Caesar} because Brutus and Antony give us a sharp public example. Brutus can make his case sound orderly. Antony then challenges whether Brutus's key premise is true.

A valid argument with a false or unproved premise is not sound. That distinction matters in Rome, in essays, in politics, and in ordinary judgment.

\focusbox{\textbf{Whately's Lesson 6 habit:} Give the argument a form test and a fact test. Do not call an argument fully successful until it passes both.}

\section*{Central Question}
If Brutus's argument has a valid form, what must still be proved before Rome should accept his conclusion?

\section*{Guided Text: Brutus in Act 3, Scene 2}
Brutus has helped kill Caesar. He now explains himself to the citizens. Read the original text below. Watch for the premise that would have to be true before his conclusion could be sound.

\speaker{Brutus}
\begin{verse}
If there be any in this assembly, any dear friend of\\
Caesar's, to him I say, that Brutus' love to Caesar\\
was no less than his. If then that friend demand\\
why Brutus rose against Caesar, this is my answer:\\
--Not that I loved Caesar less, but that I loved\\
Rome more. Had you rather Caesar were living and\\
die all slaves, than that Caesar were dead, to live\\
all free men?\gloss{Brutus makes Rome choose between Caesar alive with slavery and Caesar dead with freedom.} As Caesar loved me, I weep for him;\\
as he was fortunate, I rejoice at it; as he was\\
valiant, I honour him: but, as he was ambitious, I\\
slew him. There is tears for his love; joy for his\\
fortune; honour for his valour; and death for his\\
ambition.\gloss{The key factual premise is that Caesar's ambition threatened Rome.} Who is here so base that would be a\\
bondman? If any, speak; for him have I offended.\\
Who is here so rude that would not be a Roman? If\\
any, speak; for him have I offended. Who is here so\\
vile that will not love his country? If any, speak;\\
for him have I offended. I pause for a reply.
\end{verse}

\section*{Guided Analysis}
Brutus's argument can be reconstructed in a valid form:
\begin{itemize}[leftmargin=*]
\item \textbf{Major premise:} All men whose ambition would enslave Rome deserve death for Rome's freedom.
\item \textbf{Minor premise:} Caesar was a man whose ambition would enslave Rome.
\item \textbf{Conclusion:} Therefore, Caesar deserved death for Rome's freedom.
\end{itemize}

\subsection*{Form test}
The form is valid. If both premises are true, the conclusion follows. Brutus has given the crowd an argument that can be arranged into disciplined shape.

\subsection*{Fact test}
The pressure point is the minor premise. Has Brutus proved that Caesar's ambition would enslave Rome? He asserts it. He makes the alternative frightening. But assertion and fear are not the same as proof.

\focusbox{\textbf{Lesson 6 takeaway:} Brutus's argument may pass the form test, but it still needs the fact test. If Caesar's ambition would enslave Rome, the argument becomes much stronger. If that premise is false or unproved, the argument is not sound.}

\section*{Discussion}
\begin{enumerate}[leftmargin=*]
\item What is Brutus's conclusion?
\item Which premise is the pressure point?
\item Why might the crowd accept the premise quickly?
\item What kind of evidence would Brutus need in order to prove it?
\item Does a valid form make Brutus innocent? Why not?
\end{enumerate}

\section*{Independent Text: Antony in Act 3, Scene 2}
Antony does not first attack Brutus's form. He attacks the factual claim that Caesar was dangerously ambitious. Read the original text, then diagnose how Antony pressures Brutus's premise.

\speaker{Antony}
\begin{verse}
The noble Brutus\\
Hath told you Caesar was ambitious:\\
If it were so, it was a grievous fault,\\
And grievously hath Caesar answer'd it.\\
Here, under leave of Brutus and the rest--\\
For Brutus is an honourable man;\\
So are they all, all honourable men--\\
Come I to speak in Caesar's funeral.\\
He was my friend, faithful and just to me:\\
But Brutus says he was ambitious;\\
And Brutus is an honourable man.\\
He hath brought many captives home to Rome\\
Whose ransoms did the general coffers fill:\\
Did this in Caesar seem ambitious?\gloss{Evidence 1: Caesar enriched Rome rather than himself.}\\
When that the poor have cried, Caesar hath wept:\\
Ambition should be made of sterner stuff:\\
Yet Brutus says he was ambitious;\\
And Brutus is an honourable man.\gloss{Evidence 2: compassion does not fit Antony's picture of hard ambition.}\\
You all did see that on the Lupercal\\
I thrice presented him a kingly crown,\\
Which he did thrice refuse: was this ambition?\gloss{Evidence 3: refusing the crown challenges the claim of kingly ambition.}\\
Yet Brutus says he was ambitious;\\
And, sure, he is an honourable man.\\
I speak not to disprove what Brutus spoke,\\
But here I am to speak what I do know.
\end{verse}

\section*{Independent Analysis}
Answer in complete thoughts.
\begin{enumerate}[leftmargin=*]
\item Which premise in Brutus's argument does Antony attack?\studentline
\item List Antony's three pieces of evidence.\studentline\studentline\studentline
\item Does Antony prove Caesar was not ambitious, or does he make Brutus's premise doubtful? Explain.\studentline\studentline
\item Give Brutus's argument a final diagnosis: sound, unsound, or not yet proved. Defend your answer.\studentline\studentline
\item How does Antony use irony around the phrase ``honourable man'' to affect the fact test?\studentline\studentline
\end{enumerate}

\section*{Portfolio Artifact: Julius Caesar Form Test / Fact Test Diagnosis}
Complete the diagnosis below carefully enough that it can be saved in your Logic Portfolio.

\textbf{Argument diagnosed:}\studentline
\textbf{Conclusion:}\studentline
\textbf{Form test: valid or invalid?}\studentline
\textbf{Fact test: which premise is true, false, doubtful, or unproved?}\studentline
\textbf{Evidence for or against the pressure-point premise:}\studentline\studentline
\textbf{Final diagnosis: sound, unsound, or not yet proved?}\studentline
\textbf{One sentence explaining why validity alone is not enough:}\studentline\studentline

\section*{Closing Synthesis}
Whately's distinction protects readers from two opposite mistakes. We should not reject Brutus's argument merely because we dislike the conclusion; first we should test its form fairly. But we also should not accept Brutus's argument merely because its form can be made valid. Antony's speech shows why the fact test matters. Public judgment becomes dangerous when a crowd treats a frightening premise as proved before the evidence has been weighed.

\end{document}
